\section{Grounding, entailment, and NLI}\label{section:background-grounding-entailment-and-nli}

Checking whether a claim is supported by the evidence text it cites can be formulated as a natural-language-inference (NLI) problem: does the evidence paragraph entail the claim? In this thesis, this NLI problem is answered with a frozen biomedical NLI cross-encoder model, based on PubMedBERT~\parencite{gu2021biomedbert} and fine-tuned on the MNLI~\parencite{williams2018mnli} and MedNLI~\parencite{romanov2018mednli} entailment datasets, which build on the broader NLI tradition represented by SNLI~\parencite{}. This NLI model is applied following the sentence-pair scoring setup used in related representation-learning work~\parencite{reimers2019sbert}. 

The same NLI model is used in two different roles in the thesis pipeline: as a grounding filter, and as a relation classifier. The grounding filter scores each (claim, evidence text) pair and drops findings that are not supported enough by their evidence; the relation classifier scores the entailment and contradiction of each rule pair, and assigns a supporting, contradicting, or unrelated label to each of these pairs; the logic used for the label assignment is described in Section~\ref{section:Knowledge_Extraction}. The use of entailment scores as a faithfulness check connects the thesis to other work on detecting unsupported or hallucinated content by comparing generated claims against evidence~\parencite{}. 

The contribution of this thesis is the empirical results obtained from using an off-the-shelf biomedical NLI model with frozen thresholds, showing that such use of these models performs sufficiently well for both source faithfulness and relation classification. The obtained empirical results for both roles of the NLI model are reported in Chapter~\ref{chapter:Results}. 